% Part:sets-functions-relations
% Chapter: size-of-sets
% Section: non-enumerability

\documentclass[../../../include/open-logic-section]{subfiles}

\begin{document}

\olfileid{sfr}{siz}{nen}

\olsection{\printtoken{S}{nonenumerable} سټونه}

\begin{editorial}
  دا برخه په \olref[enm]{sec} کښې د ورکړل شوي تعريف په کارولو د
  $\Bin^\omega$ او $\Pow{\PosInt}$ د شمېر نۀ وړ والے ثابتوي۔ دا
  د \olref[enm-alt]{sec} له بڼې څخه لږه ابتدايي او لږه زياته
  تفصيلي جوړه شوې ده۔
\end{editorial}

ځينې سټونه، لکه د مثبتو صحيح عددونو سټ~$\PosInt$، نامتناهي دي۔
تر اوسه مو د نامتناهي سټونو داسې بېلګې ليدلې چې ټولې
!!{enumerable} وې۔ خو داسې نامتناهي سټونه هم شته چې دا خاصيت نۀ
لري۔ داسې سټونه \emph{!!{nonenumerable}} بلل کېږي۔

تر هر څه لومړے ښايي همدا خبره لا حيرانوونکې وي چې !!{nonenumerable}
سټونه شته۔ د هر !!{enumerable} سټ~$A$ دپاره يوه !!a{surjective}
تابع $f \colon \PosInt \to A$ شته۔ که يو سټ !!{nonenumerable} وي،
داسې تابع نۀ شته۔ يعنې هېڅ هغه تابع چې د~$\PosInt$ بې‌شمېره
!!{element}s پر~$A$ نقشوي، ټول~$A$ نۀ شي پوره کولے۔ نو په~$A$
کښې د بې‌شمېره مثبتو صحيح عددونو په پرتله «زيات» !!{element}s شته۔

څنګه به ثابتوو چې يو سټ !!{nonenumerable} دے؟ بايد وښيئ چې داسې
پر هدف سټ بشپړه تابع وجود نۀ شي لرلے۔ په مساوي ډول، بايد وښيئ
چې د~$A$ غړي په يو اړخيز نامتناهي لړ کښې نۀ شي شمېرل کېدے۔ د
دې تر ټولو ښه لار دا ده چې وښيئ د~$A$ د !!{element}s هر لړ بايد
لږ تر لږه يو غړے پرېږدي؛ يا دا چې هېڅ تابع
$f\colon \PosInt \to A$ !!{surjective} نۀ شي کېدے۔ دا د کانتور
په \emph{قطري طريقه} کولے شو۔ د~$A$ د !!{element}s يو لړ، د بېلګې
په ډول $x_1$، $x_2$، \dots، چې راکړل شي، موږ د~$A$ يو بل غړے
جوړوو چې د خپل جوړښت له مخې په هېڅ ډول په هغه لړ کښې نۀ شي راتلے۔

زموږ لومړۍ بېلګه د ټولو هغو بې‌تشې نامتناهي تسلسلونو سټ
$\Bin^\omega$ دے چې له $0$ګانو او $1$ګانو جوړ دي۔

\begin{thm}
\ollabel{thm:nonenum-bin-omega}
$\Bin^\omega$~!!{nonenumerable} دے۔
\end{thm}

\begin{proof}
د تناقض دپاره فرض کړئ چې $\Bin^\omega$ !!{enumerable} دے؛ يعنې
فرض کړئ چې د~$\Bin^\omega$ د ټولو !!{element}s يو لړ $s_{1}$،
$s_{2}$، $s_{3}$، $s_{4}$، \dots{} شته۔ له دې هر يو $s_i$ پخپله
د $0$ګانو او~$1$ګانو يو نامتناهي تسلسل دے۔ د دې لړ د $i$م تسلسل
$j$م غړے به $s_i(j)$ وبولو۔ نو $i$م تسلسل~$s_i$ دا دے:
\[
s_i(1), s_i(2), s_i(3), \dots
\]

دا لړ، او په کښې د هر تسلسل~$s_i$ غړي، په يوه جدولي ترتيب کښې
اېښودلے شو:
\[
\begin{array}{c|c|c|c|c|c}
& 1 & 2 & 3 & 4 & \dots \\\hline
1 & \mathbf{s_{1}(1)} & s_{1}(2) & s_{1}(3) & s_1(4) & \dots \\\hline
2 & s_{2}(1)& \mathbf{s_{2}(2)} & s_2(3) & s_2(4) & \dots \\\hline
3 & s_{3}(1)& s_{3}(2) & \mathbf{s_3(3)} & s_3(4) & \dots \\\hline
4 & s_{4}(1)& s_{4}(2) & s_4(3) & \mathbf{s_4(4)} & \dots \\\hline
\vdots & \vdots & \vdots & \vdots & \vdots & \mathbf{\ddots}
\end{array}
\]
د څنګ پر خوا نښې په لړ کښې د تسلسلونو شمېرې $s_1$، $s_2$،
\dots راکوي؛ د پاسنۍ خوا شمېرې د هر يوه تسلسل !!{element}s په
نښه کوي۔ د بېلګې په ډول، $s_{1}(1)$ د هغه عدد نوم دے چې، که $0$
وي يا~$1$، په تسلسل~$s_{1}$ کښې لومړے !!{element} دے؛ او همداسې
نور۔

اوس د $0$ګانو او $1$ګانو يو داسې نامتناهي تسلسل
$\overline{s}$ جوړوو چې په هېڅ ډول په دې لړ کښې نۀ شي راتلے۔ د
$\overline{s}$ تعريف به په لړ $s_1$، $s_2$، \dots پورې تړلے وي۔
د $0$ګانو او $1$ګانو د نامتناهي تسلسلونو هر نامتناهي لړ يو داسې
نامتناهي تسلسل~$\overline{s}$ راپيدا کوي چې په ډاډه توګه په هغه
لړ کښې نۀ راځي۔

د $\overline{s}$ د تعريفولو دپاره د هغه ټول !!{element}s ټاکو؛
يعنې د ټولو $n \in \PosInt$ دپاره $\overline{s}(n)$ ټاکو۔ دا کار
د پورته جدولي ترتيب پر قطر د لاندې تللو (ځکه ورته «قطري طريقه»
وايو)، بيا د هر $1$ په $0$ او د هر $0$ په~$1$ بدلولو کوو۔ په لا
مجرد ډول، $\overline{s}(n)$ د دې له مخې $0$ يا $1$ تعريفوو چې د
قطر $n$م !!{element}، $s_n(n)$، $1$ دے که $0$:
\[
\overline{s}(n) =
\begin{cases}
1 & \text{که $s_{n}(n) = 0$ وي}\\
0 & \text{که $s_{n}(n) = 1$ وي}.
\end{cases}
\]
که فارمولونه مو د حالتونو تر تعريف زيات خوښېږي، دا هم تعريفولے
شئ چې $\overline{s}(n) = 1 - s_n(n)$۔

په څرګند ډول $\overline{s}$ د $0$ګانو او $1$ګانو يو نامتناهي
تسلسل دے، ځکه دا يوازې د هغو $0$ګانو او $1$ګانو مخالف تسلسل دے
چې زموږ د جدولي ترتيب پر قطر راځي۔ نو $\overline{s}$ د
~$\Bin^\omega$ يو !!a{element} دے۔ خو دا په لړ $s_1$، $s_2$،
\dots{} کښې نۀ شي راتلے۔ ولې؟

دا په لړ کښې لومړے تسلسل~$s_1$ نۀ شي کېدے، ځکه په لومړي
!!{element} کښې له $s_1$ څخه توپير لري۔ $s_1(1)$ چې هر څه وي،
موږ $\overline{s}(1)$ د هغه مخالف تعريف کړے دے۔ دا په لړ کښې
دويم تسلسل نۀ شي کېدے، ځکه $\overline{s}$ په دويم غړي کښې له
$s_2$ څخه توپير لري: که $s_2(2)$، $0$ وي، $\overline{s}(2)$،
$1$ دے، او برعکس۔ او همداسې نور۔

په لا دقيقه توګه: که $\overline{s}$ په لړ کښې وے، يو $k$ به
داسې وے چې $\overline{s} = s_{k}$۔ دوه تسلسلونه هله او يوازې
هله يو شان وي چې په هر مقام کښې سره سمون ولري؛ يعنې د هر~$n$
دپاره $\overline{s}(n) = s_{k}(n)$ وي۔ نو په ځانګړي ډول چې
$n = k$ واخلو، بايد $\overline{s}(k) = s_{k}(k)$ وي۔ $s_k(k)$
يا $0$ دے يا~$1$۔ که $0$ وي، $\overline{s}(k)$ بايد~$1$ وي؛
همداسې مو $\overline{s}$ تعريف کړے دے۔ خو که $s_k(k) = 1$ وي،
بيا هم د $\overline{s}$ د تعريف له مخې $\overline{s}(k) = 0$۔
په هر حالت کښې $\overline{s}(k) \neq s_{k}(k)$۔

موږ د~$\Bin^\omega$ د !!{element}s د يوه لړ $s_1$، $s_2$،
\dots{} په فرضولو پيل وکړ۔ له دې لړ څخه مو تسلسل~$\overline{s}$
جوړ کړ چې ثابته مو کړه په لړ کښې نۀ شي راتلے۔ خو که ټول $s_i$
د $0$ګانو او $1$ګانو تسلسلونه وي، دا خامخا د $0$ګانو او
$1$ګانو يو تسلسل دے؛ يعنې $\overline{s} \in \Bin^\omega$۔ دا په
ځانګړي ډول ښيي چې د~$\Bin^\omega$ د \emph{ټولو} !!{element}s هېڅ
لړ نۀ شي کېدے، ځکه د هر داسې لړ دپاره يو داسې تسلسل
~$\overline{s}$ هم جوړولے شو چې په ډاډه توګه په لړ کښې نۀ راځي؛
نو دا فرض چې په~$\Bin^\omega$ کښې د ټولو تسلسلونو يو لړ شته،
تناقض ته رسېږي۔
\end{proof}

\begin{explain}
د ثبوت دې طريقې ته «قطرول» ويل کېږي، ځکه د~$\overline{s}$ د
تعريف دپاره د جدولي ترتيب قطر کاروي۔ په قطرولو کښې د جدول شتون
لازم نۀ دے: په ورته مفکوره ښودلے شو چې سټونه !!{enumerable} نۀ
دي، ان هغه وخت چې هېڅ جدول او رښتينے قطر نۀ وي۔
\end{explain}

\begin{thm}
\ollabel{thm:nonenum-pownat}
$\Pow{\PosInt}$ !!{enumerable} نۀ دے۔
\end{thm}

\begin{proof}
په هماغه طريقه مخکښې ځو: ښيو چې د~$\PosInt$ د فرعي سټونو د هر
لړ دپاره د $\PosInt$ يو داسې فرعي سټ شته چې په هغه لړ کښې نۀ
شي راتلے۔ فرض کړئ دا د~$\PosInt$ د فرعي سټونو يو ورکړل شوے لړ دے:
\[
Z_{1}, Z_{2}, Z_{3}, \dots
\]
اوس يو سټ $\overline{Z}$ داسې تعريفوو چې د هر $n \in \PosInt$
دپاره، $n \in \overline{Z}$ هله او يوازې هله چې $n \notin Z_{n}$:
\[
\overline{Z} = \Setabs{n \in \PosInt}{n \notin Z_n}
\]
$\overline{Z}$ په څرګند ډول د مثبتو صحيح عددونو يو سټ دے، ځکه
د فرض له مخې هر~$Z_n$ داسې دے؛ او له دې امله
$\overline{Z} \in \Pow{\PosInt}$۔ خو $\overline{Z}$ په لړ کښې نۀ
شي راتلے۔ د دې د ښودلو دپاره به ثابته کړو چې د هر
$k \in \PosInt$ دپاره $\overline{Z} \neq Z_k$۔

نو $k \in \PosInt$ خپله سرے ونيسئ۔ موږ $\overline{Z}$ داسې تعريف
کړے چې د هر $n \in \PosInt$ دپاره، $n \in \overline{Z}$ هله او
يوازې هله چې $n \notin Z_n$۔ په ځانګړي ډول چې $n=k$ واخلو،
$k \in \overline{Z}$ هله او يوازې هله چې $k \notin Z_k$۔ خو دا
ښيي چې $\overline{Z} \neq Z_k$، ځکه $k$ د يوه سټ !!a{element}
دے خو د بل نۀ دے، او له دې امله $\overline{Z}$ او $Z_k$ بېل
!!{element}s لري۔ ځکه چې $k$ خپله سرے و، $\overline{Z}$ په لړ
$Z_1$، $Z_2$، \dots کښې نۀ دے۔
\end{proof}

\begin{explain}
تېر ثبوت د قطر نوم وانۀ خيست، خو که داسې ئې انځور کړئ، قطري
طريقه په کښې ليدلے شئ: وانګېړئ چې سټونه $Z_1$، $Z_2$، \dots
په يوه جدولي ترتيب کښې ليکل شوي، چې هر !!{element}~$j \in Z_i$
په~$j$م ستون کښې ليکل کېږي۔ فرض کړئ په هغه لړ کښې لومړي څلور
سټونه $\{1,2,3,\dots\}$، $\{2, 4, 6, \dots\}$، $\{1,2,5\}$ او
$\{3,4,5,\dots\}$ دي۔ نو جدولي ترتيب به داسې پيل شي:
\[
\begin{array}{r@{}rrrrrrr}
  Z_1 = \{ & \mathbf{1}, & 2, & 3, & 4, & 5, & 6, & \dots\}\\
  Z_2 = \{ &  & \mathbf{2}, &  & 4, &  & 6, & \dots\}\\
  Z_3 = \{ & 1, & 2, &  &  & 5\phantom{,} &  & \}\\
  Z_4 = \{ &  &  & 3, & \mathbf{4}, & 5, & 6, & \dots\}\\
  \vdots & & & & & \ddots
\end{array}
\]
بيا $\overline{Z}$ هغه سټ دے چې د قطر پر لاندې تلو جوړېږي: هغه
شمېرې پرېږدو چې پر قطر راځي، او هغه $j$ګانې ور شاملوو چې د جدول
په $j$م قطار/ستون کښې تش ځاے لري۔ په پورته حالت کښې به $1$ او
$2$ پرېږدو،~$3$ به ور شاملوو،~$4$ به پرېږدو، او همداسې نور۔
\end{explain}

\begin{prob}
په قطري استدلال وښيئ چې $\Pow{\Nat}$ !!{nonenumerable} دے۔
\end{prob}

\begin{prob}\label{sfr:siz:nen:prob:f-posint}
په څرګند قطري استدلال وښيئ چې د تابعو سټ
$f \colon \PosInt \to \PosInt$ !!{nonenumerable} دے۔ يعنې وښيئ
چې که $f_1$، $f_2$، \dots د تابعو يو لړ وي او هره
$f_i\colon \PosInt \to \PosInt$ وي، نو يوه
$\overline{f}\colon \PosInt \to \PosInt$ شته چې په دې لړ کښې نۀ ده۔
\end{prob}

\end{document}
